% ── SECTION 7: CONCLUSION (~600 words) ───────────────────────────────────────

\section{Conclusion: The Bahá'í Contribution to Science-Religion Scholarship}
\label{sec:conclusion}

% Key points:
%   - The Bahá'í writings contribute something unique: a spiritually revealed
%     account that arrived, by a different method, at the same structural
%     conclusions as modern physics and process philosophy.
%   - This is the two-wings principle — not as hope or doctrine, but as fact.
%   - Opens new channels for academic Bahá'í-science dialogue.

The Bahá'í writings contribute something unique to the science-religion
conversation.
They constitute a spiritually revealed account of reality that arrived,
by a different method, at the same structural conclusions as modern physics
and process philosophy --- forty years before Whitehead named the framework
and a century before physics confirmed it.

This is the two-wings principle --- not as a hope, not as a doctrine, but as
a demonstrated fact.
Science and religion have, independently of each other, arrived at the same
description of what reality is made of: relationship, constituted by relational
affinity, sustained by what each tradition calls love.

The implications for Bahá'í scholarship are significant.
The principle of the harmony of science and religion can now be defended not
merely by assertion but by argument --- by showing, in the language of modern
physics and philosophy, that the structural architecture of the Bahá'í
metaphysics of \mahabbat is identical to the structural architecture that
physics has empirically confirmed at the most fundamental level of physical
description we possess.

That is not a small finding.
It is the two wings, confirmed.
